\subsection{Hệ hỗ trợ quyết định (DSS)}\label{sec:dss}

\hspace{1cm}
Theo Turban et al., một Hệ hỗ trợ Quyết định (DSS) hỗ trợ người ra quyết định trải qua các
giai đoạn: \textit{Intelligence} (nhận diện vấn đề) $\to$ \textit{Design} (xây dựng mô hình
và các phương án) $\to$ \textit{Choice} (lựa chọn phương án) $\to$ \textit{Implementation}
(triển khai). Trong bài tập lớn này:
\begin{itemize}
    \item Giai đoạn \textit{Intelligence} được hỗ trợ bởi kho dữ liệu và các OLAP view.
    \item Giai đoạn \textit{Design} được hỗ trợ bởi hai mô hình dự đoán: rủi ro giao trễ và rủi
          ro nhận đánh giá thấp. Mỗi mô hình chấm điểm từng đơn và xếp mức High/Medium/Low; với đơn
          đang vận chuyển, chức năng mô phỏng hành động (\textit{Simulate action}, bản chất là phân tích
          \textit{what-if}) cho biết rủi ro thay đổi thế nào nếu kéo dài ngày hẹn giao.
    \item Giai đoạn \textit{Choice} và \textit{Implementation} thuộc về người dùng Ops: chọn
          đơn cần can thiệp (liên hệ người bán, đổi hãng vận chuyển, báo trước khách) và chọn
          ngưỡng rủi ro để gắn khách hàng vào chiến dịch chăm sóc trên Ops Console.
\end{itemize}
Hai mô hình là \textit{model-driven DSS} kết hợp \textit{data-driven DSS}: đầu vào là dữ liệu
đã tích hợp trong kho (\texttt{FACT\_ORDER\_FULFILLMENT} cùng các dimension), đầu ra được ghi
ngược lại vào kho để dashboard và chiến dịch sử dụng.

\subsubsection{Thời điểm ra quyết định trong vòng đời đơn hàng}\label{sec:dss-timing}
\hspace{1cm}
Hình~\ref{fig:lifecycle} minh họa hai thời điểm được chọn
(tránh rò rỉ dữ liệu tương lai -- \textit{data
leakage}). 
\begin{figure}[H]
\centering
\resizebox{\textwidth}{!}{%
\begin{tikzpicture}[
    ev/.style={circle, fill=blue!60, inner sep=2.2pt},
    lbl/.style={font=\small, align=center},
    chk/.style={draw, thick, rounded corners, fill=orange!15, font=\small, align=center,
                text width=6.4cm}
]
\draw[-{Stealth[length=2.5mm]}, thick] (0,0) -- (14.6,0);
\node[ev] (p) at (0.4,0) {};   \node[lbl, below=3pt] at (p) {Mua};
\node[ev] (a) at (2.6,0) {};   \node[lbl, below=3pt] at (a) {Duyệt};
\node[ev] (c) at (5.4,0) {};   \node[lbl, below=3pt] at (c) {Bàn giao\\vận chuyển};
\node[ev] (d) at (9.4,0) {};   \node[lbl, below=3pt] at (d) {Giao tới\\khách};
\node[ev] (r) at (13.4,0) {};  \node[lbl, below=3pt] at (r) {Khách viết\\đánh giá};
\draw[dashed, gray] (11.3,-0.3) -- (11.3,0.3);
\node[lbl, text=gray, below=3pt] at (11.3,0) {Ngày hẹn};
\node[chk] (k1) at (3.9,1.6) {\textbf{Mô hình 1 -- rủi ro giao trễ}\\chấm điểm khi bàn giao VC};
\node[chk] (k2) at (10.9,1.6) {\textbf{Mô hình 2 -- rủi ro đánh giá thấp}\\chấm điểm khi vừa giao};
\draw[-{Stealth[length=2mm]}, thick, orange!70!black] (k1.south) -- (c);
\draw[-{Stealth[length=2mm]}, thick, orange!70!black] (k2.south) -- (d);
\end{tikzpicture}}
\caption{Hai thời điểm chấm điểm trong vòng đời một đơn hàng}
\label{fig:lifecycle}
\end{figure}

\begin{itemize}
    \item \textbf{Rủi ro giao trễ} được chấm khi người bán vừa bàn giao hàng cho đơn vị vận
          chuyển: lúc này đã biết hàng đi mất bao lâu và còn bao nhiêu ngày tới hạn, nhưng
          chưa biết ngày giao thực tế. Thực nghiệm cho thấy nếu chấm ngay lúc đặt hàng thì
          ROC-AUC chỉ đạt 0,693 (Mục~6), không đủ để ra quyết định.
    \item \textbf{Rủi ro đánh giá thấp} được chấm khi đơn vừa giao tới khách, trước khi khách
          viết đánh giá: đã biết đơn trễ bao nhiêu ngày -- yếu tố ảnh hưởng mạnh nhất tới điểm
          đánh giá -- và vẫn còn thời gian để chăm sóc khách.
\end{itemize}

\subsubsection{Đặc trưng (features) và nhãn}
\hspace{1cm}
Cả hai mô hình đọc cùng một bảng dữ liệu cấp đơn được tạo từ \texttt{FACT\_ORDER\_\allowbreak FULFILLMENT}
nối với \texttt{DIM\_LOCATION}, \texttt{DIM\_SELLER}, \texttt{DIM\_PRODUCT},
\texttt{DIM\_PAYMENT\_TYPE}, \texttt{DIM\_DATE} (cài đặt tại \texttt{dss/risk/features.py}).

\begin{table}[H]
\centering
\caption{Nhãn và đặc trưng của hai mô hình}
\label{tab:features}
\begin{tabular}{@{}>{\raggedright\arraybackslash}p{3.1cm}>{\raggedright\arraybackslash}p{5.3cm}>{\raggedright\arraybackslash}p{5.6cm}@{}}
\toprule
& \textbf{Rủi ro giao trễ} & \textbf{Rủi ro đánh giá thấp} \\
\midrule
Nhãn & \texttt{is\_late}: ngày giao thực tế $>$ ngày hẹn & \texttt{is\_low\_review}: điểm
đánh giá $\leq$ 2 \\
Đặc trưng chung & \multicolumn{2}{>{\raggedright\arraybackslash}p{10.9cm}}{khoảng cách người
bán -- khách (\texttt{distance\_km}), cùng bang hay không, số ngày hẹn giao, hạn bàn giao của
người bán, số giờ duyệt đơn, tháng/thứ mua, số sản phẩm, số người bán, tổng giá, phí vận chuyển
và tỷ lệ phí/giá, khối lượng, thể tích, số kỳ trả góp; biến phân loại: bang khách, bang người
bán, ngành hàng, phương thức thanh toán} \\
Đặc trưng riêng & \texttt{carrier\_days} (số ngày tới lúc bàn giao), \texttt{slack\_days} (số
ngày hẹn còn lại sau khi bàn giao), \texttt{handoff\_vs\_\allowbreak limit\_days} (người bán bàn giao muộn
so với hạn) & \texttt{carrier\_days}, \texttt{actual\_days} (tổng số ngày giao),
\texttt{delay\_days} (số ngày trễ/sớm so với hẹn), cờ trễ \\
Tập dữ liệu & đơn đã giao và đã bàn giao vận chuyển & đơn đã giao và có đánh giá \\
\bottomrule
\end{tabular}
\end{table}

\subsubsection{Thuật toán, huấn luyện và đánh giá}
\hspace{1cm}
\textbf{Thuật toán}: mô hình chính là \textbf{HistGradientBoostingClassifier} (gradient
boosting trên cây quyết định của scikit-learn): xử lý trực tiếp giá trị thiếu và biến phân
loại, nắm bắt được quan hệ phi tuyến (ví dụ: chỉ khi \texttt{slack\_days} xuống dưới một mức
nào đó thì rủi ro mới tăng vọt). Mô hình \textbf{Logistic Regression} (one-hot + chuẩn hóa)
được dùng làm \textit{baseline} để chứng minh mô hình chính tốt hơn (cài đặt tại
\texttt{dss/risk/train.py}).

\textbf{Chia dữ liệu theo thời gian} (\textit{time-based split}), không xáo trộn ngẫu nhiên,
để mô phỏng việc dùng mô hình cho đơn hàng tương lai: huấn luyện trên đơn mua từ 01/2017 đến
05/2018 (validation 03--05/2018), kiểm định trên đơn mua 06--08/2018. Dữ liệu năm 2016 quá
thưa và từ 09/2018 gần như không còn đơn nên bị loại.

\textbf{Chỉ số đánh giá}: nhãn lệch mạnh (trên tập kiểm định chỉ 3,6\% đơn trễ và 9,7\% đơn bị
đánh giá thấp), nên \textit{accuracy} không có ý nghĩa (đoán ``không trễ'' cho mọi đơn đã đạt
96,4\%). Báo cáo dùng ROC-AUC, PR-AUC (\textit{average precision}), precision/recall/F1 tại
ngưỡng quyết định, và precision/recall khi Ops chỉ xử lý 10\% số đơn có điểm cao nhất.
Mức độ quan trọng của đặc trưng được đo bằng \textit{permutation importance} (mức giảm PR-AUC
khi xáo trộn đặc trưng) trên 6.000 đơn kiểm định.

\subsubsection{Ngưỡng quyết định theo chi phí}
\hspace{1cm}
Mô hình trả về xác suất $p$; việc biến xác suất thành hành động là một quyết định nghiệp vụ.
Gọi $C_{FP}$ là chi phí một lần can thiệp không cần thiết và $C_{FN}$ là chi phí bỏ sót một
đơn có sự cố. Can thiệp có lợi khi $p \cdot C_{FN} > (1-p) \cdot C_{FP}$, tức:
\[
    p \;\geq\; t^{*} = \frac{C_{FP}}{C_{FP} + C_{FN}}
\]
\begin{itemize}
    \item \textbf{Giao trễ}: chi phí liên hệ người bán/hãng vận chuyển và báo trước cho khách
          được xem tương đương chi phí để một đơn trễ mà không báo trước $\Rightarrow$
          $t^{*} = 1/(1+1) = 0{,}5$.
    \item \textbf{Đánh giá thấp}: voucher khoảng R\$10, mất một khách do trải nghiệm xấu ước
          tính R\$40 $\Rightarrow$ $t^{*} = 10/(10+40) = 0{,}2$.
\end{itemize}
Đơn có $p \geq t^{*}$ xếp mức \textbf{High}, $t^{*}/2 \leq p < t^{*}$ là \textbf{Medium}, còn
lại là \textbf{Low}. Các tham số chi phí đặt trong \texttt{DECISION\_COSTS} và có thể điều chỉnh
theo thực tế doanh nghiệp. Nhóm cũng thử chọn ngưỡng tối đa F1 trên tập validation, nhưng cách
này thất bại với mô hình giao trễ do tỷ lệ trễ thay đổi mạnh theo thời gian (Mục~6.4).

\subsubsection{Mô phỏng hành động (what-if)}\label{sec:dss-whatif}
\hspace{1cm}
Trong bảng đơn đang vận chuyển có rủi ro giao trễ trên Ops Console, mỗi dòng có nút
\textit{Simulate action}.
Chức năng này được thêm vào để đánh giá rủi ro giao trễ cho 1 đơn nếu thay đổi ngày hẹn
(kéo dài ngày hẹn giao thêm $N$ ngày, tức tăng \texttt{estimated\_days}). Sử dụng chính mô hình đã huấn luyện
(API \texttt{POST /risk/late\_delivery/what-if})
Với đơn mà người bán bàn giao quá muộn (như đơn đầu tiên ở Hình~\ref{fig:ui-ops-late}, bàn giao sau 41 ngày), kéo dài 5 ngày gần
như không làm giảm rủi ro (98,5\% $\to$ 98,4\%), cho thấy cần hành động khác như liên hệ người bán, tự động huỷ đơn
hoặc báo trước cho khách. Giao diện chỉ có phương án kéo dài ngày hẹn vì đơn đã bàn giao cho đơn
vị vận chuyển, nên ngày hẹn là đại lượng Ops còn thay đổi được.
Mô hình đánh giá thấp không có bước này trên giao diện: đơn đã giao tới khách thì không thể thay
đổi thời gian giao, nên hành động phù hợp là chăm sóc khách hàng.

Kết quả mô hình đánh giá thấp được dùng trực tiếp để gắn khách hàng vào chiến dịch chăm sóc:
Ops chọn ngưỡng rủi ro và một chiến dịch trong danh sách tự quản lý, hệ thống lấy khách của các
đơn có điểm $\geq$ ngưỡng và ghi vào \texttt{dw.campaign\_target} (mỗi khách một lần cho mỗi
chiến dịch, xét theo khách thật \texttt{customer\_id} nên khách đã đổi địa chỉ, tức có nhiều phiên
bản SCD2, cũng không bị gắn trùng; nếu khách có nhiều đơn đạt ngưỡng thì giữ đơn rủi ro nhất). 
Vì Olist là dữ liệu lịch sử, các đơn hiện được chấm điểm thuộc tập kiểm định và đã
có review thật, nên đây là mô phỏng quy trình: với ngưỡng 50\%, hệ thống chọn 376 trong 18.520
đơn. Khi triển khai thật, danh sách sẽ lấy từ các đơn vừa giao để Ops chăm sóc trước khi khách
đánh giá.

\textbf{Công cụ dự đoán cho đơn tự nhập.} Ngoài đơn có sẵn trong kho, người dùng có thể nhập một
đơn giả định để biết rủi ro giao trễ và rủi ro đánh giá thấp
(\texttt{POST /risk/\{model\}/predict}, màn hình \texttt{\#/predict}). Mô hình cần khoảng 22
đặc trưng trong khi người dùng chỉ nhập vài thông tin chính, nên các đặc trưng còn thiếu được
điền bằng \textit{profile} lưu cùng mô hình lúc huấn luyện (trung vị trên các đơn đã giao của
giai đoạn train), theo thứ tự ưu tiên: giá trị người dùng nhập $\to$ trung vị của tuyến hoặc
ngành hàng (nếu nhóm có ít nhất 20 đơn) $\to$ trung vị toàn sàn (Bảng~\ref{tab:predict-input}).

% \begin{table}[H]
% \centering
% \caption{Đầu vào của công cụ dự đoán và giá trị mặc định khi bỏ trống}
% \label{tab:predict-input}
% \small
% \begin{tabular}{@{}>{\raggedright\arraybackslash}p{4.2cm}>{\raggedright\arraybackslash}p{3.6cm}>{\raggedright\arraybackslash}p{6.2cm}@{}}
% \toprule
% \textbf{Đầu vào} & \textbf{Dùng cho} & \textbf{Khi bỏ trống} \\
% \midrule
% Bang người bán, bang khách, ngành hàng, ngày mua, ngày hẹn giao & Cả hai mô hình & Bắt buộc
% (chỉ nhận giá trị có trong tập train) \\
% Khoảng cách (km) & Cả hai mô hình & Tự điền ngay khi chọn hai bang: trung vị của tuyến; tuyến
% dưới 20 đơn dùng khoảng cách giữa tâm hai bang \\
% Hạn bàn giao (số ngày sau khi mua) & Giao trễ & Cấu hình của sàn
% (\texttt{DEFAULT\_SHIPPING\_LIMIT\_DAYS}); không cấu hình thì trung vị dữ liệu (6 ngày) \\
% Ngày bàn giao vận chuyển & Giao trễ & Trung vị số ngày bàn giao của tuyến \\
% Ngày giao tới khách & Đánh giá thấp & Trung vị số ngày giao của tuyến \\
% Giá trị đơn, phí ship, số sản phẩm, thanh toán & Cả hai mô hình & Trung vị của ngành hàng;
% thanh toán phổ biến nhất (thẻ tín dụng) \\
% \bottomrule
% \end{tabular}
% \end{table}

Đặc trưng suy ra (\texttt{slack\_days}, \texttt{handoff\_vs\_limit\_days},
\texttt{delay\_days}) được tính lại từ các ngày bằng cùng hàm \texttt{derive()} như lúc huấn
luyện. Kết quả trả về kèm nguồn của từng giá trị (nhập, trung vị tuyến, trung vị ngành, quy định
sàn, \ldots) và tỷ lệ lịch sử của tuyến để người dùng biết đâu là giả định. 
Ví dụ tuyến
SP $\to$ AL (220 đơn lịch sử, 27,3\% trễ), hẹn giao 27 ngày: nếu người bán bàn giao đúng mức
điển hình, rủi ro trễ khoảng 50\% (vừa chạm ngưỡng High); nếu bàn giao vào ngày thứ 19, rủi ro tăng
lên 87\% (High) (Hình~\ref{fig:ui-predict}).

Mô hình học \textit{tương quan} trong dữ liệu lịch sử chứ không phải quan hệ nhân quả.
Thay đổi những đại lượng ít biến động trong dữ liệu (ví dụ đặt hạn bàn giao 3 ngày, trong khi
dữ liệu chủ yếu 5--7 ngày) có thể cho kết quả ngược trực giác. Vì vậy công cụ phù hợp nhất để so
sánh các kịch bản về ngày bàn giao và ngày giao -- hai đặc trưng quan trọng nhất của hai mô hình
(Mục~6.3).

% \subsubsection{Kiến trúc phục vụ mô hình}
\hspace{1cm}
Hai mô hình được đóng gói chung trong một dịch vụ FastAPI (\texttt{dss/api/main.py}). Mỗi lần
huấn luyện lại (\texttt{POST /risk/retrain}, do DAG \texttt{dss\_batch} gọi) sẽ ghi kết quả vào
kho dữ liệu:
\begin{itemize}
    \item \texttt{dw.model\_metrics} -- chỉ số đánh giá của mô hình chính và baseline;
    \item \texttt{dw.model\_feature\_importance} -- permutation importance;
    \item \texttt{dw.order\_risk\_score} -- điểm và mức rủi ro của từng đơn kiểm định
          (\texttt{scope = holdout}, kèm kết quả thật để đối chiếu) và của 1.106 đơn đang vận
          chuyển (\texttt{scope = open}).
\end{itemize}
% Các endpoint chính: \texttt{GET /risk/\{model\}/summary}, \texttt{GET /risk/\{model\}/orders},
% \texttt{GET /risk/metrics}, \texttt{GET /risk/\{model\}/importance},
% \texttt{POST /risk/\{model\}/what-if}, \texttt{POST /risk/\{model\}/predict},
% \texttt{GET /risk/predict/distance}. Mã nguồn chia theo vai trò: \texttt{risk/features.py} (đọc dữ
% liệu, đặc trưng, nhãn), \texttt{risk/train.py} (huấn luyện, đánh giá, chấm điểm và ghi kết quả vào kho),
% \texttt{risk/predict.py} (dùng mô hình đã huấn luyện cho mô phỏng hành động (what-if) và dự đoán đơn giả định; file mô hình
% được giữ trong bộ nhớ và tự đọc lại khi huấn luyện lại). Profile dùng cho công cụ dự đoán được lưu cùng mô hình trong
% file \texttt{.pkl}, nên mỗi lần huấn luyện lại thì các giá trị điển hình cũng được cập nhật theo
% dữ liệu mới. Việc tách \texttt{dss} thành dịch vụ độc lập với
% \texttt{backend-api} cho phép huấn luyện lại mô hình (thông qua DAG \texttt{dss\_batch}) độc
% lập với vòng đời của tầng API phục vụ frontend, đồng thời dễ mở rộng hoặc thay thế từng mô
% hình mà không ảnh hưởng các thành phần khác.

\subsection{Luồng dữ liệu đầu-cuối}
\hspace{1cm}
Hình~\ref{fig:sequence} minh họa luồng xử lý một yêu cầu điển hình từ Manager Dashboard: (1)
người dùng chọn khoảng thời gian trên dashboard; (2) frontend gọi API
\texttt{GET /dashboard/\allowbreak revenue-trend}; (3) backend-api truy vấn \texttt{mv\_revenue\_by\_month\_category}
trên kho dữ liệu; (4) kết quả trả về và được trực quan hóa dưới dạng biểu đồ.

\begin{figure}[H]
\centering
\begin{tikzpicture}[
    actor/.style={draw, rounded corners, fill=black!5, minimum width=2.8cm,
                  minimum height=1.2cm, font=\small, align=center},
    lbl/.style={font=\scriptsize, align=center},
    arr/.style={-{Stealth[length=2mm]}, thick}
]
\node[actor] (fe) at (0,0) {Manager\\Dashboard};
\node[actor] (be) at (7,0) {Backend API\\(BFF)};
\node[actor] (dw) at (14,0) {Kho dữ liệu\\(OLAP view)};

\draw[arr] ([yshift=5pt]fe.east) -- node[above, lbl]{\texttt{GET /dashboard/revenue-trend}\\\texttt{?year=\&category=\&region=}}
    ([yshift=5pt]be.west);
\draw[arr] ([yshift=5pt]be.east) -- node[above, lbl]{\texttt{SELECT ... FROM}\\\texttt{mv\_revenue\_by\_}\\\texttt{month\_category}}
    ([yshift=5pt]dw.west);
\draw[arr] ([yshift=-8pt]dw.west) -- node[below, lbl]{Các dòng kết quả}
    ([yshift=-8pt]be.east);
\draw[arr] ([yshift=-8pt]be.west) -- node[below, lbl]{JSON}
    ([yshift=-8pt]fe.east);
\end{tikzpicture}
\caption{Luồng xử lý một yêu cầu truy vấn doanh thu từ dashboard}
\label{fig:sequence}
\end{figure}

\texttt{backend-api} (FastAPI) đóng vai trò \textit{Backend-for-Frontend}: các truy vấn thống kê
(KPI, doanh thu, hiệu suất giao hàng, xếp hạng sản phẩm) đọc trực tiếp từ OLAP view và bảng
fact của kho dữ liệu; các yêu cầu liên quan tới mô hình (danh sách đơn rủi ro, mô phỏng hành động, dự đoán
cho đơn tự nhập) được chuyển tiếp sang \texttt{dss} service; thao tác chiến dịch chăm sóc (thêm,
sửa, xóa, gắn khách có rủi ro đánh giá thấp $\geq$ ngưỡng) ghi vào \texttt{dw.care\_campaign} và
\texttt{dw.campaign\_target}.
